


\begin{Activite}[][5 min]
Un joueur de foot tire un coup franc. La hauteur du ballon (exprimée en mètres) par rapport au sol en fonction du temps $t$ (exprimé en secondes) est donnée par la fonction $h$ définie sur l'intervalle $[0; 6]$ par :
$$h(t) = -t^2 + 6t$$

\begin{enumerate}
    \item Calculer $h(0)$ et $h(6)$. Que signifient ces deux résultats pour le ballon ?
    \item Calculer la hauteur du ballon après $3$ secondes.
    \item 
    \begin{enumerate}
    \item Montrer que pour tout $t \in [0; 6]$, on a : $h(t) - 9 = -(t - 3)^2$
    \item En déduire que pour tout $t \in [0; 6]$, on a : $h(t) \leq 9$
    \end{enumerate}
    \item Quelle est la hauteur maximale atteinte par le ballon ? À quel instant est-elle atteinte ?
\end{enumerate}

\end{Activite}